\part{Écritures, simplifications}

\section{Simplification sous forme d'une fractions}\label{convfrac}

\subsection{Idée}

\begin{tipblock}
L'idée est d'obtenir une commande pour \textit{simplifier} un calcul sous forme de fraction irréductible.
\end{tipblock}

\begin{PresCodeTexPL}{listing only}
\ConversionFraction(*)[option de formatage]{calcul}
\end{PresCodeTexPL}

\subsection{Commande et options}

\begin{cautionblock}
Quelques explications sur cette commande :

\begin{itemize}
	\item \cmaj{2.5.1} la version \textit{étoilée} force l'écriture du signe \og $-$ \fg{} sur le numérateur ;
	\item le premier argument, \textit{optionnel} et entre \textsf{[...]} permet de spécifier un formatage du résultat :
	\begin{itemize}
			\item \Cle{t} pour l'affichage de la fraction en mode \textsf{tfrac} ;
			\item \Cle{d} pour l'affichage de la fraction en mode \textsf{dfrac} ;
			\item \Cle{n} pour l'affichage de la fraction en mode \textsf{nicefrac} ;
			\item \Cle{dec} pour l'affichage du résultat en mode \texttt{décimal} (sans arrondi !) ;
			\item \Cle{dec=k} pour l'affichage du résultat en mode \texttt{décimal} arrondi à $10^{-k}$ ;
		\end{itemize}
	\item le second argument, \textit{obligatoire}, est quant à lui, le calcul en syntaxe \textsf{xint}.
\end{itemize}

À noter que la macro est dans un bloc \ctex{ensuremath} donc les \ctex{\$...\$} ne sont pas nécessaires.
\end{cautionblock}

\begin{PresCodeTexPL}{listing only}
\ConversionFraction{-10+1/3*(-5/16)}          %sortie par défaut 
\ConversionFraction*{-10+1/3*(-5/16)}         %sortie fraction avec - sur numérateur
\ConversionFraction[d]{-10+1/3*(-5/16)}       %sortie en displaystyle
\ConversionFraction[n]{-10+1/3*(-5/16)}       %sortie en nicefrac
\ConversionFraction[dec=4]{-10+1/3*(-5/16)}   %sortie en décimal arrondi à 0,0001
\ConversionFraction{2+91/7}                   %entier formaté
\ConversionFraction{111/2145}
\ConversionFraction{111/3}
\end{PresCodeTexPL}

\begin{PresCodeSortiePL}{text only}
\ConversionFraction{-10+1/3*(-5/16)}

\smallskip

\ConversionFraction*{-10+1/3*(-5/16)}

\smallskip

\ConversionFraction[d]{-10+1/3*(-5/16)}

\smallskip

\ConversionFraction[n]{-10+1/3*(-5/16)}

\smallskip

\ConversionFraction[dec=4]{-10+1/3*(-5/16)}

\smallskip

\ConversionFraction{2+91/7}

\smallskip

\ConversionFraction{111/2145}

\smallskip

\ConversionFraction{111/3}
\end{PresCodeSortiePL}

\begin{PresCodePL}{}
$\frac{111}{2145}=\ConversionFraction{111/2145}$ \\

$\frac{3}{15}=\ConversionFraction[]{3/15}$ \\

$\tfrac{3}{15}=\ConversionFraction[t]{3/15}$ \\

$\dfrac{3}{15}=\ConversionFraction[d]{3/15}$ \\

$\dfrac{0,42}{0,015}=\ConversionFraction[d]{0.42/0.015}$ \\

$\dfrac{0,41}{0,015}=\ConversionFraction[d]{0.41/0.015}$ \\

$\dfrac{1}{7}-\dfrac{3}{8}=\ConversionFraction[d]{1/7-3/8}$ \\

$\ConversionFraction[d]{1+1/2}$ \\

$\ConversionFraction{0.1/0.7+30/80}$
\end{PresCodePL}

\begin{noteblock}
A priori le package \ctex{xint} permet de s'en sortir pour des calculs \og simples \fg, je ne garantis pas que tout calcul ou toute division donne un résultat \textit{satisfaisant} !
\end{noteblock}

\newpage

\section{Écriture d'un trinôme, trinôme aléatoire}\label{trinome}

\subsection{Idée}

\begin{tipblock}
L'idée est de proposer une commande pour écrire, sous forme développée réduite, un trinôme en fonction de ses coefficients $a$, $b$ et $c$ (avec $a\neq0$), avec la gestion des coefficients nuls ou égaux à $\pm1$.

\smallskip

En combinant avec le package \ctex{xfp} et fonction de générateur d'entiers aléatoires, on peut de ce fait proposer une commande pour générer aléatoirement des trinômes à coefficients entiers (pour des fiches d'exercices par exemple).

\smallskip

L'affichage des monômes est géré par le package \ctex{siunitx} et le tout est dans un environnement \ctex{ensuremath}.
\end{tipblock}

\begin{PresCodeTexPL}{listing only}
\EcritureTrinome[options]{a}{b}{c}
\end{PresCodeTexPL}

\begin{PresCodePL}{}
\EcritureTrinome{1}{7}{0}\\
\EcritureTrinome{1.5}{7.3}{2.56}\\
\EcritureTrinome{-1}{0}{12}\\
\EcritureTrinome{-1}{-5}{0}
\end{PresCodePL}

\subsection{Clés et options}

\begin{cautionblock}
Quelques clés et options sont disponibles :

\begin{itemize}
	\item la clé booléenne \Cle{Alea} pour autoriser les coefficients aléatoires ;\hfill{}défaut \Cle{false}
	\item la clé booléenne \Cle{Anegatif} pour autoriser $a$ à être négatif.\hfill{}défaut \Cle{true}
\end{itemize}
\vspace*{-\baselineskip}\leavevmode
\end{cautionblock}

\begin{noteblock}
La clé \Cle{Alea} va modifier la manière de saisir les coefficients, il suffira dans ce cas de  préciser les bornes, sous la forme \ctex{valmin,valmax}, de chacun des coefficients. C'est ensuite le package \ctex{xfp} qui va se charger de générer les coefficients.
\end{noteblock}

\begin{PresCodePL}{}
Avec $a$ entre 1 et 5 (et signe aléatoire) puis $b$ entre $-5$ et 5 puis $c$ entre $-10$ et 20 : \\
$f(x)=\EcritureTrinome[Alea]{1,5}{-5,5}{-10,10}$\\
$g(x)=\EcritureTrinome[Alea]{1,5}{-5,5}{-10,10}$\\
$h(x)=\EcritureTrinome[Alea]{1,5}{-5,5}{-10,10}$\\
Avec $a$ entre 1 et 10 (forcément positif) puis $b$ entre $-2$ et 2 puis $c$ entre 0 et 4 : \\
\EcritureTrinome[Alea,Anegatif=false]{1,10}{-2,2}{0,4}\\
\EcritureTrinome[Alea,Anegatif=false]{1,10}{-2,2}{0,4}\\
\EcritureTrinome[Alea,Anegatif=false]{1,10}{-2,2}{0,4}
\end{PresCodePL}

\newpage

\section{Simplification de racines}\label{simplracine}

\subsection{Idée}

\begin{tipblock}
\cmaj{2.1.0} L'idée est de proposer une commande pour simplifier \textit{automatiquement} une racine carrée, sous la forme $\frac{a\sqrt{b}}{c}$ avec $\frac{a}{c}$ irréductible et $b$ le \frquote{plus petit possible}.
\end{tipblock}

\begin{PresCodeTexPL}{listing only}
\SimplificationRacine{expression ou calcul}
\end{PresCodeTexPL}

\begin{PresCodePL}{}
\SimplificationRacine{48} \\ \SimplificationRacine{100/34}\\
\SimplificationRacine{99999} \\ \SimplificationRacine{1500*0.31*(1-0.31)}\\
\end{PresCodePL}

\begin{noteblock}
C'est -- comme souvent -- le package \ctex{xint} qui s'occupe en interne des calculs, et qui devrait donner des résultats satisfaisants dans la majorité des cas (attention aux \textit{grands nombres}\ldots)

\smallskip

La commande ne fait pas office de \textit{calculatrice}, elle ne permet \textit{que} de simplifier \textit{une} racine carrée (donc transformer si besoin !).
\end{noteblock}

\subsection{Exemples}

\begin{PresCodePL}{}
%Simplification d'un module de complexe
$\left| 4+6\text{i}\right| = \sqrt{4^2+6^2} = \sqrt{\xinteval{4**2+6**2}}=\SimplificationRacine{4**2+6**2}$

%Simplification n°1
$\frac{1}{\sqrt{6}}=\left(\sqrt{\frac{1}{6}}\right)=\SimplificationRacine{1/6}$

%Simplification n°2
$\frac{42}{\sqrt{5}}=\left(\sqrt{\frac{42^2}{5}}\right)=\SimplificationRacine{(42*42)/5}$

%Écart-type d'une loi binomiale
$\sqrt{\num{150}\times\num{0.35}\times(1-\num{0.35})} = \displaystyle\SimplificationRacine{150*0.35*(1-0.35)}$
\end{PresCodePL}

\newpage

\section{Mesure principale d'un angle}\label{mesureprincipale}

\subsection{Idée}

\begin{tipblock}
\cmaj{2.1.2} L'idée est de proposer (sur une suggestion de Marylyne Vignal) une commande pour déterminer la mesure principale d'un angle en radian.
\end{tipblock}

\begin{PresCodeTexPL}{listing only}
\MesurePrincipale[booléens]{angle}   %dans un mode mathématique
\end{PresCodeTexPL}

\begin{noteblock}
La commande est à insérer dans un environnement mathématique, via \ctex{\$...\$} ou \ctex{\textbackslash[...\textbackslash]}.

L'angle est donné sous forme \textit{explicite} avec la chaîne \ctex{pi}.
\end{noteblock}

\subsection{Exemples}

\begin{cautionblock}
Pour cette commande :

\begin{itemize}
	\item le booléen \Cle{d} permet de forcer l'affichage en \ctex{displaystyle} ; \hfill{}défaut \Cle{false}
	\item le booléen \Cle{Crochets} permet d'afficher le \textit{modulo} entre crochets (sinon parenthèses) ;
	
	\hfill{}défaut \Cle{false}
	\item \cmaj{2.6.0} le booléen \Cle{Brut} pour afficher uniquement la mesure principale ; \hfill{}défaut \Cle{false}
	\item l'argument \textit{obligatoire} est en écriture \textit{en ligne}.
\end{itemize}
\vspace*{-\baselineskip}\leavevmode
\end{cautionblock}

\begin{PresCodeTexPL}{listing only}
$\MesurePrincipale[d]{54pi/7}$
$\MesurePrincipale[d]{-128pi/15}$
$\MesurePrincipale{3pi/2}$
$\MesurePrincipale[Crochets]{5pi/2}$
$\MesurePrincipale{-13pi}$
$\MesurePrincipale{28pi}$
$\MesurePrincipale[d]{14pi/4}$
$\MesurePrincipale[Crochets]{14pi/7}$
$\dfrac{121\pi}{12} = \MesurePrincipale[Brut]{121pi/12}$ à $2\pi$ près
\end{PresCodeTexPL}

\begin{PresCodeSortiePL}{text only}
$\MesurePrincipale[d]{54pi/7}$

\medskip

$\MesurePrincipale{-128pi/15}$

\medskip

$\MesurePrincipale{3pi/2}$

\medskip

$\MesurePrincipale[Crochets]{5pi/2}$

\medskip

$\MesurePrincipale{-13pi}$

\medskip

$\MesurePrincipale{28pi}$

\medskip

$\MesurePrincipale[d]{14pi/4}$

\medskip

$\MesurePrincipale[Crochets]{14pi/7}$

\medskip

$\dfrac{121\pi}{12} = \MesurePrincipale[d,Brut]{121pi/12}$ à $2\pi$ près
\end{PresCodeSortiePL}

\pagebreak

\section{Lignes trigonométriques}\label{lignestrigo}

\subsection{Idée}

\begin{tipblock}
\cmaj{2.6.0} L'idée est de proposer pour déterminer les lignes trigonométriques (cos, sin et tan) d'angles classiques, formés des \og $\pi$ \fg{} et \og $\pi$ sur 2 ; 3 ; 4 ; 5 ; 6 ; 8 ; 10 ; 12 \fg{}.

\smallskip

La commande détermine -- et affiche si demandée la réduction -- et la valeur exacte de la ligne trigonométrique demandée.
\end{tipblock}

\begin{PresCodeTexPL}{listing only}
\LigneTrigo(*)[booléens]{cos/sin/tan}(angle)
\end{PresCodeTexPL}

\subsection{Commande}

\begin{cautionblock}
Pour cette commande :

\begin{itemize}
	\item la version \textit{étoilée} n'affiche pas l'angle initial ;
	\item le booléen \Cle{d} permet de forcer l'affichage en \ctex{displaystyle} ; \hfill{}défaut \Cle{false}
	\item le booléen \Cle{Etapes} permet d'afficher la réduction avant le résultat ; \hfill{}défaut \Cle{false}
	\item le premier argument \textit{obligatoire}, entre \ctex{\{...\}} est le type de calcul demandé, parmi \Cle{cos / sin / tan} ;
	\item le second argument \textit{obligatoire}, entre \ctex{(...)} est l'angle, donné en ligne, avec \ctex{pi}.
\end{itemize}
\vspace*{-\baselineskip}\leavevmode
\end{cautionblock}

\begin{PresCodePL}{}
$\LigneTrigo{cos}(56pi/3)$ et $\LigneTrigo{sin}(56pi/3)$ et $\LigneTrigo{tan}(56pi/3)$
\end{PresCodePL}

\begin{PresCodePL}{}
$\LigneTrigo[d,Etapes]{cos}(56pi/3)$ et $\LigneTrigo[d,Etapes]{sin}(56pi/3)$
\end{PresCodePL}

\begin{PresCodePL}{}
$\LigneTrigo*[d,Etapes]{cos}(2pi/3)$ et $\LigneTrigo*[d,Etapes]{sin}(2pi/3)$
\end{PresCodePL}

\begin{PresCodePL}{}
$\LigneTrigo[d,Etapes]{cos}(146pi)$ et $\LigneTrigo[d,Etapes]{sin}(146pi)$
\end{PresCodePL}

\begin{PresCodePL}{}
$\LigneTrigo[d,Etapes]{cos}(-551pi/12)$ et $\LigneTrigo[d,Etapes]{sin}(-551pi/12)$
\end{PresCodePL}

\begin{PresCodePL}{}
$\LigneTrigo[d,Etapes]{cos}(447pi/8)$ et $\LigneTrigo[d,Etapes]{sin}(447pi/8)$
\end{PresCodePL}

\begin{PresCodePL}{}
$\LigneTrigo*[d,Etapes]{cos}(-pi/8)$ et $\LigneTrigo*[d,Etapes]{sin}(-pi/8)$
\end{PresCodePL}

\begin{PresCodePL}{}
$\LigneTrigo[d,Etapes]{cos}(-595pi/12)$ et $\LigneTrigo[d,Etapes]{sin}(-595pi/12)$ et $\LigneTrigo[d,Etapes]{tan}(-595pi/12)$
\end{PresCodePL}

\begin{PresCodePL}{}
$\LigneTrigo[d,Etapes]{cos}(33pi/10)$ et $\LigneTrigo[d,Etapes]{sin}(33pi/10)$\\
$\LigneTrigo[d,Etapes]{tan}(33pi/10)$
\end{PresCodePL}

\begin{PresCodePL}{}
$\LigneTrigo[d,Etapes]{cos}(-14pi/5)$ et $\LigneTrigo[d,Etapes]{sin}(-14pi/5)$\\
$\LigneTrigo[d,Etapes]{tan}(-14pi/5)$
\end{PresCodePL}

\subsection{Valeurs disponibles}

\begin{noteblock}
Les valeurs disponibles sont :

\begin{tblr}{hlines,vlines,colspec={Q[1cm,m,l]*{9}{Q[1.195cm,m,c]}},cells={font=\scriptsize},row{1}={bg=lightgray!50},column{1}={bg=lightgray!50}}
	angle	& $0$	& $\nicefrac{\pi}{6}$	& $\nicefrac{\pi}{4}$	& $\nicefrac{\pi}{3}$	& $\nicefrac{\pi}{2}$	& $\nicefrac{2\pi}{3}$	& $\nicefrac{3\pi}{4}$	& $\nicefrac{5\pi}{6}$ & $\pi$ \\
	cos		& $\LigneTrigo{cos}(0)$ & $\LigneTrigo{cos}(pi/6)$ & $\LigneTrigo{cos}(pi/4)$ & $\LigneTrigo{cos}(pi/3)$ & $\LigneTrigo{cos}(pi/2)$ & $\LigneTrigo{cos}(2pi/3)$ & $\LigneTrigo{cos}(3pi/4)$ & $\LigneTrigo{cos}(5pi/6)$ & $\LigneTrigo{cos}(pi)$ \\
	sin		& $\LigneTrigo{sin}(0)$ & $\LigneTrigo{sin}(pi/6)$ & $\LigneTrigo{sin}(pi/4)$ & $\LigneTrigo{sin}(pi/3)$ & $\LigneTrigo{sin}(pi/2)$ & $\LigneTrigo{sin}(2pi/3)$ & $\LigneTrigo{sin}(3pi/4)$ & $\LigneTrigo{sin}(5pi/6)$ & $\LigneTrigo{sin}(pi)$ \\
	tan		& $\LigneTrigo{tan}(0)$ & $\LigneTrigo{tan}(pi/6)$ & $\LigneTrigo{tan}(pi/4)$ & $\LigneTrigo{tan}(pi/3)$ & $\LigneTrigo{tan}(pi/2)$ & $\LigneTrigo{tan}(2pi/3)$ & $\LigneTrigo{tan}(3pi/4)$ & $\LigneTrigo{tan}(5pi/6)$ & $\LigneTrigo{tan}(pi)$ \\
\end{tblr}

\medskip

\begin{tblr}{hlines,vlines,colspec={Q[1cm,m,l]*{9}{Q[1.195cm,m,c]}},cells={font=\scriptsize},row{1}={bg=lightgray!50},column{1}={bg=lightgray!50}}
	angle	& 	& $\nicefrac{-\pi}{6}$	& $\nicefrac{-\pi}{4}$	& $\nicefrac{-\pi}{3}$	& $\nicefrac{-\pi}{2}$	& $\nicefrac{-2\pi}{3}$	& $\nicefrac{-3\pi}{4}$	& $\nicefrac{-5\pi}{6}$ &  \\
	cos		&  & $\LigneTrigo{cos}(-pi/6)$ & $\LigneTrigo{cos}(-pi/4)$ & $\LigneTrigo{cos}(-pi/3)$ & $\LigneTrigo{cos}(-pi/2)$ & $\LigneTrigo{cos}(-2pi/3)$ & $\LigneTrigo{cos}(-3pi/4)$ & $\LigneTrigo{cos}(-5pi/6)$ &  \\
	sin		&  & $\LigneTrigo{sin}(-pi/6)$ & $\LigneTrigo{sin}(-pi/4)$ & $\LigneTrigo{sin}(-pi/3)$ & $\LigneTrigo{sin}(-pi/2)$ & $\LigneTrigo{sin}(-2pi/3)$ & $\LigneTrigo{sin}(-3pi/4)$ & $\LigneTrigo{sin}(-5pi/6)$ &  \\
	tan		&  & $\LigneTrigo{tan}(-pi/6)$ & $\LigneTrigo{tan}(-pi/4)$ & $\LigneTrigo{tan}(-pi/3)$ & $\LigneTrigo{tan}(-pi/2)$ & $\LigneTrigo{tan}(-2pi/3)$ & $\LigneTrigo{tan}(-3pi/4)$ & $\LigneTrigo{tan}(-5pi/6)$ &  \\
\end{tblr}

\medskip

\begin{tblr}{hlines,vlines,colspec={Q[1cm,m,l]*{8}{Q[1.4cm,m,c]}},cells={font=\scriptsize},row{1}={bg=lightgray!50},column{1}={bg=lightgray!50}}
	angle	& $\nicefrac{\pi}{8}$	& $\nicefrac{3\pi}{8}$	& $\nicefrac{5\pi}{8}$	& $\nicefrac{7\pi}{8}$	& $\nicefrac{\pi}{12}$	& $\nicefrac{5\pi}{12}$	& $\nicefrac{7\pi}{12}$ & $\nicefrac{11\pi}{12}$  \\
	cos		& $\LigneTrigo{cos}(pi/8)$ & $\LigneTrigo{cos}(3pi/8)$ & $\LigneTrigo{cos}(5pi/8)$ & $\LigneTrigo{cos}(7pi/8)$ & $\LigneTrigo{cos}(pi/12)$ & $\LigneTrigo{cos}(5pi/12)$ & $\LigneTrigo{cos}(7pi/12)$ & $\LigneTrigo{cos}(11pi/12)$  \\
	sin		& $\LigneTrigo{sin}(pi/8)$ & $\LigneTrigo{sin}(3pi/8)$ & $\LigneTrigo{sin}(5pi/8)$ & $\LigneTrigo{sin}(7pi/8)$ & $\LigneTrigo{sin}(pi/12)$ & $\LigneTrigo{sin}(5pi/12)$ & $\LigneTrigo{sin}(7pi/12)$ & $\LigneTrigo{sin}(11pi/12)$  \\
	tan		& $\LigneTrigo{tan}(pi/8)$ & $\LigneTrigo{tan}(3pi/8)$ & $\LigneTrigo{tan}(5pi/8)$ & $\LigneTrigo{tan}(7pi/8)$ & $\LigneTrigo{tan}(pi/12)$ & $\LigneTrigo{tan}(5pi/12)$ & $\LigneTrigo{tan}(7pi/12)$ & $\LigneTrigo{tan}(11pi/12)$  \\
\end{tblr}

\medskip

\begin{tblr}{hlines,vlines,colspec={Q[1cm,m,l]*{8}{Q[1.4cm,m,c]}},cells={font=\scriptsize},row{1}={bg=lightgray!50},column{1}={bg=lightgray!50}}
	angle	& $\nicefrac{-\pi}{8}$	& $\nicefrac{-3\pi}{8}$	& $\nicefrac{-5\pi}{8}$	& $\nicefrac{-7\pi}{8}$	& $\nicefrac{-\pi}{12}$	& $\nicefrac{-5\pi}{12}$	& $\nicefrac{-7\pi}{12}$ & $\nicefrac{-11\pi}{12}$  \\
	cos		& $\LigneTrigo{cos}(-pi/8)$ & $\LigneTrigo{cos}(-3pi/8)$ & $\LigneTrigo{cos}(-5pi/8)$ & $\LigneTrigo{cos}(-7pi/8)$ & $\LigneTrigo{cos}(-pi/12)$ & $\LigneTrigo{cos}(-5pi/12)$ & $\LigneTrigo{cos}(-7pi/12)$ & $\LigneTrigo{cos}(-11pi/12)$  \\
	sin		& $\LigneTrigo{sin}(-pi/8)$ & $\LigneTrigo{sin}(-3pi/8)$ & $\LigneTrigo{sin}(-5pi/8)$ & $\LigneTrigo{sin}(-7pi/8)$ & $\LigneTrigo{sin}(-pi/12)$ & $\LigneTrigo{sin}(-5pi/12)$ & $\LigneTrigo{sin}(-7pi/12)$ & $\LigneTrigo{sin}(-11pi/12)$  \\
	tan		& $\LigneTrigo{tan}(-pi/8)$ & $\LigneTrigo{tan}(-3pi/8)$ & $\LigneTrigo{tan}(-5pi/8)$ & $\LigneTrigo{tan}(-7pi/8)$ & $\LigneTrigo{tan}(-pi/12)$ & $\LigneTrigo{tan}(-5pi/12)$ & $\LigneTrigo{tan}(-7pi/12)$ & $\LigneTrigo{tan}(-11pi/12)$  \\
\end{tblr}

\medskip

\begin{tblr}{hlines,vlines,colspec={Q[1cm,m,l]*{8}{Q[1.4cm,m,c]}},cells={font=\scriptsize},row{1}={bg=lightgray!50},column{1}={bg=lightgray!50}}
	angle	& $\nicefrac{-4\pi}{5}$	& $\nicefrac{-3\pi}{5}$	& $\nicefrac{-2\pi}{5}$	& $\nicefrac{-\pi}{5}$	& $\nicefrac{\pi}{5}$	& $\nicefrac{2\pi}{5}$	& $\nicefrac{3\pi}{5}$ & $\nicefrac{4\pi}{5}$  \\
	cos		& $\LigneTrigo{cos}(-4pi/5)$ & $\LigneTrigo{cos}(-3pi/5)$ & $\LigneTrigo{cos}(-2pi/5)$ & $\LigneTrigo{cos}(-pi/5)$ & $\LigneTrigo{cos}(pi/5)$ & $\LigneTrigo{cos}(2pi/5)$ & $\LigneTrigo{cos}(3pi/5)$ & $\LigneTrigo{cos}(4pi/5)$  \\
	sin		& $\LigneTrigo{sin}(-4pi/5)$ & $\LigneTrigo{sin}(-3pi/5)$ & $\LigneTrigo{sin}(-2pi/5)$ & $\LigneTrigo{sin}(-pi/5)$ & $\LigneTrigo{sin}(pi/5)$ & $\LigneTrigo{sin}(2pi/5)$ & $\LigneTrigo{sin}(3pi/5)$ & $\LigneTrigo{sin}(4pi/5)$  \\
	tan		& $\LigneTrigo{tan}(-4pi/5)$ & $\LigneTrigo{tan}(-3pi/5)$ & $\LigneTrigo{tan}(-2pi/5)$ & $\LigneTrigo{tan}(-pi/5)$ & $\LigneTrigo{tan}(pi/5)$ & $\LigneTrigo{tan}(2pi/5)$ & $\LigneTrigo{tan}(3pi/5)$ & $\LigneTrigo{tan}(4pi/5)$  \\
\end{tblr}

\medskip

\begin{tblr}{hlines,vlines,colspec={Q[1cm,m,l]*{8}{Q[1.4cm,m,c]}},cells={font=\scriptsize},row{1}={bg=lightgray!50},column{1}={bg=lightgray!50}}
	angle	& $\nicefrac{-9\pi}{10}$	& $\nicefrac{-7\pi}{10}$	& $\nicefrac{-3\pi}{10}$	& $\nicefrac{-\pi}{10}$	& $\nicefrac{\pi}{10}$	& $\nicefrac{3\pi}{10}$	& $\nicefrac{7\pi}{10}$ & $\nicefrac{9\pi}{10}$  \\
	cos		& $\LigneTrigo{cos}(-9pi/10)$ & $\LigneTrigo{cos}(-7pi/10)$ & $\LigneTrigo{cos}(-3pi/10)$ & $\LigneTrigo{cos}(-pi/10)$ & $\LigneTrigo{cos}(pi/10)$ & $\LigneTrigo{cos}(3pi/10)$ & $\LigneTrigo{cos}(7pi/10)$ & $\LigneTrigo{cos}(9pi/10)$  \\
	sin		& $\LigneTrigo{sin}(-9pi/10)$ & $\LigneTrigo{sin}(-7pi/10)$ & $\LigneTrigo{sin}(-3pi/10)$ & $\LigneTrigo{sin}(-pi/10)$ & $\LigneTrigo{sin}(pi/10)$ & $\LigneTrigo{sin}(3pi/10)$ & $\LigneTrigo{sin}(7pi/10)$ & $\LigneTrigo{sin}(9pi/10)$  \\
	tan		& $\LigneTrigo{tan}(-9pi/10)$ & $\LigneTrigo{tan}(-7pi/10)$ & $\LigneTrigo{tan}(-3pi/10)$ & $\LigneTrigo{tan}(-pi/10)$ & $\LigneTrigo{tan}(pi/10)$ & $\LigneTrigo{tan}(3pi/10)$ & $\LigneTrigo{tan}(7pi/10)$ & $\LigneTrigo{tan}(9pi/10)$  \\
\end{tblr}
\end{noteblock}

\newpage

\section{Écriture sous forme de fraction irréductible d'un décimal périodique}\label{fracperiod}

\subsection{Idées}

\begin{tipblock}
\cmaj{3.00f} L'idée est de proposer une commande pour travailler sur l'écriture fractionnaire d'un nombre à écriture décimale périodique.

\smallskip

À partir de l'écriture sous la forme \og $A.B\overline{C}$ \fg, le code se charge de :

\begin{itemize}
	\item faire les calculs pour obtenir le numérateur et le dénominateur de la fraction (sans simplification) ;
	\item ou de présenter le résultat de manière basique ;
	\item ou de rédiger complètement la résolution.
\end{itemize}
\vspace*{-\baselineskip}\leavevmode
\end{tipblock}

\begin{PresCodeTexPL}{listing only}
%commande générale
\FractionPeriode[clés]{partie avant la période}{période}
%résultats bruts, pour utilisation 'externe'
\FractionPeriode{partie avant la période}{période}
%présentation simple := nombre + fraction brute + frction simplifiée
\FractionPeriode[Simple]{partie avant la période}{période}
%présentation complète := résolution
\FractionPeriode[Solution]{partie avant la période}{période}
\end{PresCodeTexPL}

\subsection{Options et clés}

\begin{cautionblock}
Les clés disponibles pour cette commande sont :

\begin{itemize}
	\item le booléen \Cle{Brut} qui permet de stocker le numérateur et dénominateur :
	\begin{itemize}
		\item dans les variables \ctex{\textbackslash FracPerNum} et \ctex{\textbackslash FracPerDenom} pour la version \textit{complète} ;
		\item dans les variables \ctex{\textbackslash FracPerNumSimpl} et \ctex{\textbackslash FracPerDenomSimpl} pour la version \textit{irréductible} ;
	\end{itemize}
	\hfill{}défaut \Cle{true}
	\item le booléen \Cle{d} pour forcer l'afficher avec \ctex{\textbackslash displaystyle} ; \hfill{}défaut \Cle{true}
	\item la clé \Cle{Inconnue} qui gère l'inconnue dans la rédaction ; \hfill{}défaut \Cle{x}
	\item le booléen \Cle{Solution} qui afficher la résolution \textit{complète} ; \hfill{}défaut \Cle{false}
	\item le booléen \Cle{Simple} qui affiche uniquement la \textit{conclusion} ; \hfill{}défaut \Cle{false}
	\item le booléen \Cle{Enonce} (en mode \Cle{Simple=true}) qui affiche le nombre de départ (formaté).
	
	\hfill{}défaut \Cle{false}
\end{itemize}

À noter que la commande gère automatiquement le mode mathématique éventuel.
\end{cautionblock}

\begin{PresCodePL}{}
%version brute, pour extraire numérateur et dénominateur puis mise en forme manuelle
\FractionPeriode{45.1}{23}Numérateur : \FracPerNum{} et \FracPerNumSimpl\par Dénominateur = \FracPerDenom{} et \FracPerDenomSimpl\par
\medskip
On a $\num{45.1}\overline{23}=\dfrac{\num{\FracPerNum}}{\num{\FracPerDenom}}=
\dfrac{\num{\FracPerNumSimpl}}{\num{\FracPerDenomSimpl}}=
\nicefrac{\num{\FracPerNumSimpl}}{\num{\FracPerDenomSimpl}}$.
\end{PresCodePL}

\begin{PresCodePL}{}
\FractionPeriode[Simple]{45.1}{23}
\end{PresCodePL}

\begin{PresCodePL}{}
\FractionPeriode[Solution]{45.1}{23}
\end{PresCodePL}

\begin{PresCodePL}{}
\FractionPeriode[Simple]{0.}{4}
\end{PresCodePL}

\begin{PresCodePL}{}
\FractionPeriode[Solution,Inconnue=n]{154.99}{174}
\end{PresCodePL}